\documentclass[12pt,a4paper]{article}
\usepackage[UTF8]{ctex}
\usepackage{amsmath,amsfonts,amssymb}
\usepackage{geometry}
\usepackage{fancyhdr}
\usepackage{enumitem}
\usepackage{xcolor}
\usepackage{tcolorbox} % 用于美化方框

% 页面设置
\geometry{left=2cm,right=2cm,top=2.5cm,bottom=2.5cm}

% 页眉页脚设置
\pagestyle{fancy}
\fancyhf{}
\fancyhead[L]{\textbf{高阶导数专题突破}}
\fancyhead[R]{自用}
\fancyfoot[C]{\thepage}
\renewcommand{\headrulewidth}{0.4pt}

% 段落设置
\setlength{\parindent}{0pt}
\setlength{\parskip}{0.8em}

% 自定义颜色和环境
\definecolor{mainblue}{RGB}{0, 102, 204}
\newtcolorbox{tiptip}{colback=blue!5!white,colframe=mainblue,title=\textbf{�� 解题思路与技巧}}

\begin{document}

\section*{专题：高阶导数}

\textbf{【解题目标】} 掌握常见函数的高阶导数公式、莱布尼兹公式、泰勒展开法及奇偶性简化技巧。

\rule{\linewidth}{0.5pt}

\section{归纳法与基本 $n$ 阶导数公式}

\textbf{【解题要点】}
求高阶导数时，若函数形式简单，推荐直接套用五大必背公式。

\begin{tcolorbox}[title=\textbf{★ 必背公式表 (建议板书)}, colback=white, colframe=black]
\begin{enumerate}[label=\arabic*)]
    \item $(e^{ax})^{(n)} = a^n e^{ax}$
    \item $[\sin(ax)]^{(n)} = a^n \sin\left(ax + \dfrac{n\pi}{2}\right)$
    \item $[\cos(ax)]^{(n)} = a^n \cos\left(ax + \dfrac{n\pi}{2}\right)$
    \item $\left(\dfrac{1}{x+a}\right)^{(n)} = \dfrac{(-1)^n n!}{(x+a)^{n+1}}$
    \item $(\ln(x+a))^{(n)} = \dfrac{(-1)^{n-1} (n-1)!}{(x+a)^n} \quad (n \geq 1)$
\end{enumerate}
\end{tcolorbox}

\vspace{1em}

% --- 例题 1 ---
\textbf{【例1】}（2023 湖南专升本 改编）
设函数 $y = \dfrac{1}{x+2}$，则 $y^{(n)}(0) = \_\_\_\_\_$。


\vspace{12em}

% --- 例题 2 ---
\textbf{【例2】}（2022 安徽专升本 真题）
函数 $y = \ln(1+x)$ 在 $x=0$ 处的三阶导数 $y^{(3)}(0) = \_\_\_\_\_$。
\vspace{12em}

\newpage

\section{莱布尼兹公式}

\textbf{【解题要点】}
适用场景：\textbf{乘积型函数}的高阶导数，即 $y=u(x) \cdot v(x)$。
\textbf{通用展开式}（类比二项式展开）：
\begin{align*}
(uv)^{(n)} &= C_n^0 u v^{(n)} + C_n^1 u' v^{(n-1)} + C_n^2 u'' v^{(n-2)} + \dots + C_n^n u^{(n)} v \\
&= u v^{(n)} + n u' v^{(n-1)} + \dfrac{n(n-1)}{2} u'' v^{(n-2)} + \dots + u^{(n)} v
\end{align*}
\textbf{破题技巧}：
\begin{itemize}
    \item 必有一项为\textbf{多项式}（如 $x, x^2$），将其设为 $u(x)$。
    \item 因为多项式求导几次后会变成 0，公式中绝大多数项为 0，只需计算前几项（\textbf{非零项数 = 多项式次数 + 1}）。
\end{itemize}

\begin{tcolorbox}[colback=green!5!white, colframe=green!50!black, title=\textbf{★ 基础补给站：组合系数 $C_n^k$ 速算}]
\textbf{1. 计算公式}：
\[ C_n^k = \frac{n!}{k!(n-k)!} = \frac{\overbrace{n \cdot (n-1) \cdots (n-k+1)}^{k\text{个}}}{\underbrace{k \cdot (k-1) \cdots 1}_{k\text{的阶乘}}} \]

\textbf{2. 常用数值（必背）}：
\begin{itemize}
    \item $C_n^0 = 1, \quad C_n^n = 1$
    \item $C_n^1 = n$
    \item $C_n^2 = \dfrac{n(n-1)}{2}$
    \item $C_n^3 = \dfrac{n(n-1)(n-2)}{6}$
\end{itemize}

\textbf{3. 对称性}：$C_n^k = C_n^{n-k}$（算大数时用），例如 $C_{10}^8 = C_{10}^2 = \frac{10 \times 9}{2} = 45$。

\textbf{4. 常考实例（直接背诵）}：
\begin{itemize}
    \item $n=2$ (二阶导)：$1, \mathbf{2}, 1$ \quad ($C_2^0, C_2^1, C_2^2$)
    \item $n=3$ (三阶导)：$1, \mathbf{3}, \mathbf{3}, 1$ \quad ($C_3^0, C_3^1, C_3^2, C_3^3$)
    \item $n=4$ (四阶导)：$1, \mathbf{4}, \mathbf{6}, \mathbf{4}, 1$ \quad ($C_4^0, C_4^1, C_4^2, C_4^3, C_4^4$)
\end{itemize}
\end{tcolorbox}

\newpage
% --- 例题 3 ---
\textbf{【例3】}（2024 河南专升本 改编）
设 $f(x) = x^2 e^x$，则 $f^{(3)}(0) = \_\_\_\_\_$。
\vspace{16em}

% --- 例题 4 ---
\textbf{【例4】}（2021 江苏专升本 真题）
函数 $f(x) = x \cdot 2^x$，则 $f''(0) = \_\_\_\_\_$。
\vspace{12em}
\newpage

\section{泰勒公式逆用求导法}

\textbf{【解题要点】}
根据麦克劳林公式（$x=0$ 处的泰勒公式）：
\[
f(x) = f(0) + f'(0)x + \frac{f''(0)}{2!}x^2 + \dots + \frac{f^{(n)}(0)}{n!}x^n + o(x^n)
\]
观察可知：\textbf{$x^n$ 的系数 $a_n$ 等于 $\frac{f^{(n)}(0)}{n!}$}。
由此得出求高阶导数的捷径：
\[
f^{(n)}(0) = n! \cdot a_n \quad (\text{即：} n! \times x^n \text{的系数})
\]

\vspace{1em}

% --- 例题 5 ---
\textbf{【例5】}（2023 广东专升本 改编）
已知 $f(x) = \dfrac{1}{1+x}$，则 $f^{(4)}(0) = \_\_\_\_\_$。
\vspace{14em}
\newline
% --- 例题 6 ---
\textbf{【例6】}（重要变式）
设 $ f(x) = \begin{cases} \dfrac{\ln(1+2x)}{x}, & x \neq 0 \\ 2, & x=0 \end{cases} \quad \left(-\dfrac{1}{2} < x < \dfrac{1}{2}\right) $，求 $ f^{(100)}(0) $。
\vspace{12em}
\newpage

\section{利用奇偶性简化计算}

\textbf{【解题要点】}
\begin{itemize}
    \item 若 $f(x)$ 是\textbf{偶函数} $\Rightarrow$ 奇数阶导数在 $x=0$ 处为 0。
    \item 若 $f(x)$ 是\textbf{奇函数} $\Rightarrow$ 偶数阶导数在 $x=0$ 处为 0。
\end{itemize}

\vspace{1em}

% --- 例题 7 ---
\textbf{【例7】}（2022 山东专升本 真题）
设 $f(x) = x^2 \cos x$，则 $f^{(3)}(0) = \_\_\_\_\_$。

\end{document}